\documentclass[UTF8,a4paper,11pt]{ctexart}

\usepackage[a4paper,left=25mm,right=25mm,top=23mm,bottom=23mm]{geometry}
\usepackage{amsmath}
\usepackage{array}
\usepackage{booktabs}
\usepackage{enumitem}
\usepackage{fancyhdr}
\usepackage{longtable}
\usepackage{listings}
\usepackage{tabularx}
\usepackage{xcolor}
\usepackage[breaklinks]{hyperref}

\definecolor{linkblue}{RGB}{30,80,150}
\definecolor{codegray}{RGB}{245,246,248}
\definecolor{rulegray}{RGB}{90,90,90}

\hypersetup{
  colorlinks=true,
  linkcolor=linkblue,
  urlcolor=linkblue,
  citecolor=linkblue,
  pdftitle={PlayTools 技术报告},
  pdfauthor={代码调查报告}
}

\setlist{nosep,leftmargin=2em}
\setlength{\parindent}{2em}
\setlength{\parskip}{0.35em}
\setlength{\emergencystretch}{6em}
\sloppy
\setcounter{tocdepth}{2}

\lstset{
  basicstyle=\ttfamily\small,
  backgroundcolor=\color{codegray},
  frame=single,
  framerule=0.3pt,
  rulecolor=\color{rulegray},
  columns=fullflexible,
  breaklines=true,
  keepspaces=true,
  showstringspaces=false,
  xleftmargin=1em,
  xrightmargin=1em
}

\newcommand{\pathref}[2]{\path{#1}:#2}
\newcommand{\localmark}{\textbf{[L]}}
\newcommand{\externalmark}{\textbf{[X]}}
\newcommand{\runtimeMark}{\textbf{[R]}}
\newcommand{\inferenceMark}{\textbf{[I]}}

\title{PlayTools 技术报告\\
  \large 注入式 iOS-on-Mac 运行时与 MaaTools 截图服务}
\author{代码调查报告}
\date{2026 年 9 月 1 日}

\begin{document}
\maketitle
\thispagestyle{empty}

\begin{abstract}
PlayTools 是 PlayCover 用于增强 iOS 应用在 macOS 上运行能力的外部 framework/dylib。本报告以 PlayCover 的 \texttt{MaaTools} 分支和其锁定的 PlayTools revision 为对象，重点说明 PlayTools 的注入方式、目标进程内的初始化链、\texttt{AKInterface} 窗口桥接、MaaTools TCP 协议及其截图像素输出。调查结论是：PlayCover 主程序负责设置、注入、签名和启动；真正的窗口采集及 TCP 服务运行在被注入的目标应用进程内，不由 PlayCover 主进程的菜单或截图 API 完成。服务返回内存中的原始像素，不保存 PNG/JPEG 文件。
\end{abstract}

\tableofcontents
\clearpage

\section{调查范围与证据边界}

\subsection{对象与版本}

本报告调查的宿主仓库是 PlayCover 的 \texttt{MaaTools} 分支，当前提交为 \texttt{63cf605}，对应版本标记 \texttt{v3.1.0.maa.7}。该分支的 \pathref{Cartfile}{1} 指向 \texttt{hguandl/PlayTools} 的 \texttt{MaaTools} 分支；\pathref{Cartfile.resolved}{1} 将依赖固定到 revision \texttt{7aa1977}。完整提交标识如下：

\begin{lstlisting}
PlayCover HEAD: 63cf605d209de834a333b09c85ca5bea261dc9eb
PlayTools revision: 7aa19776dbb72b0ae7f0989568f69c7c9c2d6543
\end{lstlisting}

PlayCover 工程的 macOS deployment target 为 12.0。当前检出目录没有 \texttt{Carthage/Build}，因此生成的 \texttt{PlayTools.xcframework} 和 \texttt{AKInterface.bundle} 不属于本地版本库的源码文件。外部 PlayTools 的静态描述仅以上述固定 revision 的源码为依据；下文会明确区分本地证据与外部证据，避免把未生成的二进制行为误写成当前本地可运行的结果。

\subsection{证据标签}

\begin{itemize}
  \item \localmark{}：PlayCover 仓库中的本地静态代码或配置。
  \item \externalmark{}：固定 revision 的外部 PlayTools 源码。
  \item \runtimeMark{}：需要运行已构建 framework、目标应用或 TCP 客户端才能验证的事实。
  \item \inferenceMark{}：根据代码结构作出的分析判断，不等同于已执行的运行时测试。
\end{itemize}

\section{总体架构与进程边界}

PlayTools 不是一个由 PlayCover 主进程通过菜单调用的独立截图工具，而是随目标应用启动、在目标应用地址空间中运行的注入式组件。其主要边界如下：

\begin{center}
\small
\begin{tabularx}{\textwidth}{>{\raggedright\arraybackslash}p{0.22\textwidth} X}
\toprule
组件或边界 & 责任 \\
\midrule
PlayCover 主程序 & 提供 SwiftUI 设置界面，保存 plist，安装和签名 PlayTools，修改目标 Mach-O，并启动目标应用。\\
目标应用进程内的 \texttt{PlayTools.dylib} & 通过 Mach-O 加载命令被注入；包含 \texttt{PlayLoader}、\texttt{PlayCover.launch} 和 MaaTools 初始化逻辑。\\
目标应用内的 \texttt{AKInterface.bundle} & 由插件加载器加载，桥接 macOS 窗口属性、窗口图像、窗口内容区域和终止操作。\\
外部 MAA 客户端 & 通过 TCP 连接目标进程内的 MaaTools，发送命令并自行解码原始像素或执行支持的输入控制。\\
\bottomrule
\end{tabularx}
\end{center}

\clearpage
\begin{samepage}
运行关系可以简化为：

\begin{lstlisting}
PlayCover host process
  +-- AppSettings / Installer / signing and injection
  `-- launch target application
       |
       v
Target iOS-on-Mac process
  +-- PlayTools.dylib
  |    `-- PlayLoader constructor
  |         `-- PlayCover.launch
  |              `-- MaaTools.initialize
  `-- AKInterface.bundle
       `-- windowImage / windowContentRect / terminateApplication
       |
       v
MAA client ---- TCP ---- MaaTools

commands: SCRN / BGR\1 / SIZE / VERN / BNDL / RECT / TOUCH / TERM
\end{lstlisting}
\end{samepage}

该边界解释了三个容易混淆的对象：iTunes 元数据中的 \texttt{screenshotUrls} 只是数据模型字段，不是截图服务；\texttt{METAL\_*} 环境变量清理属于调试环境管理，也不是 MaaTools 截图生成；PlayCover 主进程本身没有截图文件管理 API。

\section{PlayCover 侧的安装、注入与配置}

\subsection{安装决策}

安装 IPA 时，\pathref{PlayCover/PlayCover/AppInstaller/Installer.swift}{11--35} 弹窗询问是否安装 PlayTools，并将用户选择保存为安装偏好。\pathref{PlayCover/PlayCover/AppInstaller/Installer.swift}{94--98} 在导出场景调用 \texttt{PlayTools.injectInIPA}；在普通安装场景中，只有用户选择安装时才调用 \texttt{PlayTools.installInIPA}。用户选择不注入，目标应用就不会具备后文所述的运行时服务。

PlayTools 的宿主侧安装方法会把内置 framework 复制到用户的 framework 目录：\pathref{PlayCover/PlayCover/Utils/PlayTools.swift}{45--70}。针对 IPA 的安装流程包括：

\begin{enumerate}
  \item 使用注入库向目标 Mach-O 添加 PlayTools 的 dylib 加载命令；
  \item 复制本地化资源和 \texttt{AKInterface.bundle}；
  \item 清除可能导致 codesign 失败的扩展属性；
  \item 修正插件可执行权限并对插件签名；
  \item 对处理后的目标应用或其 Mach-O 重新签名。
\end{enumerate}

\pathref{PlayCover/PlayCover/Utils/PlayTools.swift}{73--125} 实现了 framework 路径下注入 dylib、复制插件和签名的主要步骤。工程构建阶段还从 \path{Carthage/Build/PlayTools.xcframework} 复制 framework，并将 \texttt{AKInterface.bundle} 复制到 framework 的插件目录：\pathref{PlayCover/PlayCover.xcodeproj/project.pbxproj}{587, 597--599}。

\subsection{目标应用启动前检查}

\pathref{PlayCover/PlayCover/Model/PlayApp.swift}{63--115} 的 \texttt{launch()} 在启动前同步设置、检查 entitlements 和 Info.plist 签名、解锁必要的 keychain 状态，并检查 PlayTools 是否已安装。若检测到目标应用已有 PlayTools，\pathref{PlayCover/PlayCover/Model/PlayApp.swift}{92--95} 会补齐插件；如果 \texttt{isInstalled()} 返回失败，则记录错误并不会进入 \texttt{runAppExec()}。因此截图服务的第一道前置条件是 PlayTools 已正确安装或注入。

\subsection{MaaTools 配置}

宿主侧的 \texttt{AppSettingsData} 定义了：

\begin{lstlisting}
maaTools     = false
maaToolsPort = 1717
\end{lstlisting}

对应代码位于 \pathref{PlayCover/PlayCover/Model/AppSettings.swift}{9--31, 104--169}。配置以 XML property list 保存到：

\begin{lstlisting}
~/Library/Containers/io.playcover.PlayCover/App Settings/<bundle-id>.plist
\end{lstlisting}

设置页的 Bypasses tab（包括 MaaTools）在没有 PlayTools 时禁用，相关逻辑位于 \pathref{PlayCover/PlayCover/Views/AppSettingsView.swift}{78--107}；MaaTools 开关和端口控件位于 \pathref{PlayCover/PlayCover/Views/AppSettingsView.swift}{539--586}，端口 UI 范围为 1024--65535。

\section{目标进程内的运行时入口}

外部 PlayTools 的 \texttt{PlayLoader.m:332--334} 使用 constructor：

\begin{lstlisting}
static void __attribute__((constructor)) initialize(void) {
    [PlayCover launch];
}
\end{lstlisting}

这意味着 dylib 一旦被目标进程加载，初始化会自动发生，不需要 PlayCover 主进程再次发送“开始截图”命令。固定 revision 中的 \texttt{PlayCover.swift:13--23} 依次初始化 \texttt{AKInterface}、PlayScreen、PlayInput 和其他运行时组件，再在主线程异步调用 \texttt{MaaTools.shared.initialize}。

MaaTools 初始化首先检查 \texttt{PlaySettings.shared.maaTools}。如果关闭则直接返回；如果开启，则每秒检查一次 key window、native scale、native bounds 和窗口标题，直到宽高非零且标题可用，然后创建 TCP listener。该等待循环没有超时，目标窗口永久不满足条件时，初始化任务也不会自然结束。

listener 就绪后，服务端将窗口标题修改为原窗口标题加上 \texttt{[localhost:<port>]}。这个标题后缀是可观察的启动副作用，但不能据此证明 socket 只绑定 loopback 地址。

\section{窗口图像采集流程}

\subsection{窗口识别与采集后端}

\texttt{AKInterface} 首先取得目标应用第一扇窗口的 \texttt{CGWindowID}。macOS 14.4 及以上使用 ScreenCaptureKit：

\begin{itemize}
  \item 以当前进程的可共享内容查找相同 window ID；
  \item 使用独立窗口过滤器和 \texttt{SCScreenshotManager.captureImage}；
  \item 配置 32-bit BGRA、sRGB、关闭光标、忽略阴影和全局裁剪，并选择最佳采集分辨率。
\end{itemize}

macOS 14.4 以下使用 \texttt{CGWindowListCreateImage}，并设置最佳分辨率、忽略窗口边框和不透明选项。需要注意，现代路径的 ScreenCaptureKit 异常只记录错误并返回空结果，不会在同一次调用中自动回退到旧 API。PlayCover 的 deployment target 是 12.0，因此两条路径都会被保留。

\subsection{标题栏裁剪与两种像素输出}

截图逻辑每次从 \texttt{AKInterface.windowImage} 获取 CGImage。它使用窗口宽高比例计算标题栏高度，并将顶部标题栏裁掉，使输出尺寸对应目标应用的内容区域。

\begin{description}
  \item[\texttt{SCRN}] 分配 \(4 \times width \times height\) 字节的 buffer，使用 sRGB CGContext 将图像绘制到 4 bytes-per-pixel 的原始缓冲区，再返回该缓冲区。代码没有进行 PNG/JPEG 编码。
  \item[\texttt{BGR\textbackslash{}1}] 通过 Accelerate/vImage 将图像的 RGBA8888 数据转换为去 alpha 的 RGB888 数据。虽然命令名含有 BGR，当前实现使用的是 \texttt{vImageConvert\_RGBA8888toRGB888}，报告应按实现记录为 RGB888 原始字节，而不是根据命令名推断通道顺序。
\end{description}

如果无法取得图像，普通截图返回长度为零的数据；BGR 响应则返回零尺寸和空数据。失败主要通过 OSLog 记录，PlayTools 没有本地文件落盘路径、命名规则或图像扩展名。

\section{MaaTools TCP 协议}

\subsection{握手与请求帧}

固定 revision 的协议版本常量为 3。新连接先读取 4 字节握手：客户端必须发送 \texttt{MAA\textbackslash{}0}，服务端回复 ASCII \texttt{OKAY}。握手之后，读取循环先读取 2 字节大端序长度，再读取相应长度的 payload。payload 的前 4 字节是命令 magic。

\begin{lstlisting}
connect:
    client  ->  MAA\0
    server  ->  OKAY

request frame:
    u16be payload_length || payload
    payload[0..3] = command magic
\end{lstlisting}

\subsection{命令与响应}

\begin{center}
\small
\begin{tabularx}{\textwidth}{>{\ttfamily\raggedright\arraybackslash}p{0.16\textwidth} X}
\toprule
命令 & 当前实现的行为 \\
\midrule
SCRN & 返回 \texttt{u32be(data\_length)}，后接每像素 4 字节的 sRGB 原始数据。\\
BGR\textbackslash{}1 & 返回 \texttt{u32be(width)}、\texttt{u32be(height)}、\texttt{u32be(data\_length)}，后接 RGB888 原始数据。\\
SIZE & 返回两个 16-bit 大端序整数：窗口内容宽度和高度。\\
VERN & 返回 32-bit 大端序协议版本，目前为 3。\\
BNDL & 返回长度和目标应用 bundle identifier 的 UTF-8 字节。\\
RECT & 返回窗口 frame 与 content rect 的坐标和尺寸字段。\\
TOUCH & 从 payload 中读取阶段与坐标，派发 down、move 或 up 触控事件。\\
TERM & 调用 \texttt{AKInterface.terminateApplication()} 终止目标应用。\\
\bottomrule
\end{tabularx}
\end{center}

截图响应中的长度字段和像素数据是协议数据，不是文件格式。客户端需要知道宽度、高度、每像素字节数和通道排列，完成内存解码后才能显示或自行保存成图片。服务端本身不提供文件名、扩展名、PNG/JPEG 编码或文件路径。

\section{配置生命周期与异常边界}

\subsection{两侧设置模型}

外部 PlayTools 的 \texttt{PlaySettings.swift:13--23} 根据当前用户名拼出与宿主相同的 plist 路径，并一次性解码 \texttt{maaTools} 与 \texttt{maaToolsPort}。文件读取或解码失败时回退到默认设置，其中 \texttt{maaTools=false}、端口为 1717，并打印失败信息。由于 MaaTools 是目标进程 constructor 启动后初始化，设置改变通常要在目标应用下一次启动时才会生效。

服务端通过 \texttt{UInt16(maaToolsPort \& 0xffff)} 生成端口；正常 UI 范围可避免大部分非法输入，但端口冲突和 listener 创建失败仍需单独处理。代码未指定 host，也没有认证或 TLS。\texttt{NWListener} 的实际绑定范围应在具备构建产物的环境中通过 socket 检查确认，不能仅凭窗口标题中的 \texttt{localhost} 判断。

\subsection{权限边界}

\pathref{PlayCover/PlayCover/Utils/Entitlements.swift}{45--65} 为目标应用生成的基础 entitlements 包含 \texttt{com.apple.security.network.server}；这是网络服务能力，与屏幕采集授权不是同一个机制。PlayCover 自身的 \texttt{PlayCover.entitlements} 和 \texttt{PlayCoverRelease.entitlements} 为空或接近为空，不代表目标应用已经获得屏幕录制许可。

\pathref{PlayCover/PlayCover/Info.plist}{78--81} 只有 \texttt{NSSystemAdministration\allowbreak{}UsageDescription}，未发现显式 Screen Recording UsageDescription、屏幕采集 entitlement 或 TCC 授权请求 API。因此 ScreenCaptureKit 最终能否成功，仍取决于 macOS 的 TCC 状态和用户实际授予的屏幕录制权限。当前代码对拒绝或采集失败主要记录错误并返回空数据。

\subsection{调试捕获变量不是 MaaTools}

\pathref{PlayCover/PlayCover/Model/PlayApp.swift}{123--150, 156--162} 在启动目标应用前清理 \texttt{DYLD\_*}、\texttt{METAL\_CAPTURE\_*}、\texttt{METAL\_FRAME\_CAPTURE\_ENABLED} 和 \texttt{MTLCaptureEnabled} 等变量。这是为了避免调试 wrapper、验证层或 Metal Frame Capture 配置从父进程继承；它不生成 MaaTools 截图，且可能使另一个独立的 Metal 调试流程失效。

\section{安全性与可靠性评估}

以下项目是基于源码的分析建议，不应理解为当前实现已经具备这些防护：

\begin{enumerate}
  \item \textbf{限制网络暴露。} listener 创建时没有显式 host、认证或加密。应明确绑定 loopback，或引入认证令牌和 TLS；同时对“未指定 host”的实际行为做运行时确认。
  \item \textbf{区分截图与控制权限。} 同一连接除了读取截图，还支持 \texttt{TOUCH} 和 \texttt{TERM}。若服务可被非预期客户端访问，可能造成输入注入或远程终止，建议将控制命令单独授权。
  \item \textbf{增加资源边界。} 请求长度字段是 16-bit，但代码未体现连接级超时、来源校验、速率限制或明确的应用层超大请求策略；应在客户端断开、半包、端口冲突和重复连接场景加入测试。
  \item \textbf{处理窗口生命周期。} 无限窗口轮询可能使初始化任务长期存在；窗口关闭、重建或应用异常退出时，listener 的取消和资源释放也应有明确策略。
  \item \textbf{改善权限反馈。} ScreenCaptureKit 失败目前主要表现为日志或空数据。宿主可以在设置或首次连接时给出 TCC 权限状态和可操作的错误信息。
  \item \textbf{固定并审计依赖。} 发布构建应继续锁定 PlayTools revision，同时保留协议版本、像素格式、权限行为和签名流程的回归测试。
\end{enumerate}

\section{分支差异与构建限制}

当前 \texttt{MaaTools} 分支不是 PlayCover 主线的普遍能力：

\begin{itemize}
  \item 当前 \pathref{Cartfile}{1} 指向 \texttt{hguandl/PlayTools} 的 \texttt{MaaTools} 分支，并在 \pathref{Cartfile.resolved}{1} 固定 revision。
  \item 对比的 \texttt{origin/develop} 仍使用 PlayTools 的另一版本，不包含当前 \texttt{maaTools} 和 \texttt{maaToolsPort} 宿主设置链。
  \item 提交 \texttt{74f55d3} 引入了 MaaTools 设置字段与 UI；因此切换到没有这些改动的分支后，不能假定截图 TCP 服务仍存在。
  \item 构建脚本依赖 Carthage 生成的 \texttt{PlayTools.xcframework}，并对 framework、插件和目标应用执行多步复制、转换及签名。当前工作树没有 \texttt{Carthage/Build}，无法仅凭宿主仓库完成真实截图服务的端到端运行验证。
  \item 仓库内未发现 XCTest、截图测试或其他专门测试目录；协议解码、TCC 拒绝、端口冲突和关闭生命周期目前缺少本地自动化覆盖。
\end{itemize}

\section{结论}

PlayTools 的核心价值不在于提供一个 PlayCover 菜单截图功能，而在于把运行时能力注入目标 iOS-on-Mac 应用。MaaTools 是其中一个可选服务：目标 dylib 被加载后自动初始化，读取宿主生成的 plist，等待 key window 就绪，调用 AKInterface 获取窗口图像，并通过 TCP 返回原始像素。PlayCover 负责让这条链路成立——安装、注入、签名、配置、目标应用启动——但截图采集与协议处理本身属于外部 PlayTools。

当前实现适合由受信任的本机客户端使用，但应谨慎对待网络边界和控制命令：协议没有内建认证，服务还支持触控和终止操作，ScreenCaptureKit 又依赖外部 TCC 授权。若要将 MaaTools 用作稳定的自动化截图接口，优先级应是确认 socket 绑定范围、增加连接认证与超时、为协议和生命周期补测试，并提供清晰的屏幕录制权限反馈。

\appendix
\section{源码证据索引}

\small
\begin{longtable}{>{\raggedright\arraybackslash}p{0.09\textwidth} >{\raggedright\arraybackslash}p{0.32\textwidth} >{\raggedright\arraybackslash}p{0.42\textwidth}}
\toprule
标签 & 文件或固定 revision & 支持的事实 \\
\midrule
\endfirsthead
\toprule
标签 & 文件或固定 revision & 支持的事实 \\
\midrule
\endhead
\bottomrule
\endfoot
\localmark & \pathref{PlayCover/PlayCover/Model/PlayApp.swift}{63--115} & 目标应用启动、PlayTools 检查、插件安装和执行入口。 \\
\localmark & \pathref{PlayCover/PlayCover/AppInstaller/Installer.swift}{11--35, 94--98} & 安装时是否注入 PlayTools 的用户选择与分支。 \\
\localmark & \pathref{PlayCover/PlayCover/Utils/PlayTools.swift}{45--70, 73--125} & framework 安装、dylib 注入、插件复制和签名。 \\
\localmark & \pathref{PlayCover/PlayCover/Model/AppSettings.swift}{9--31, 104--169} & MaaTools 默认值、端口和 plist 保存位置。 \\
\localmark & \pathref{PlayCover/PlayCover/Views/AppSettingsView.swift}{78--107, 539--586} & 无 PlayTools 时禁用设置页，以及 MaaTools 控件。 \\
\localmark & \pathref{PlayCover/PlayCover/Utils/Entitlements.swift}{45--65} & 目标应用生成的网络服务 entitlement。 \\
\localmark & \pathref{PlayCover/PlayCover/Info.plist}{78--81} & 当前 Info.plist 的权限描述范围。 \\
\localmark & \pathref{PlayCover/Cartfile.resolved}{1} & PlayTools 固定 revision。 \\
\localmark & \pathref{PlayCover/PlayCover.xcodeproj/project.pbxproj}{587, 597--599, 801, 839} & Carthage framework 复制脚本和 macOS deployment target。 \\
\externalmark & \texttt{hguandl/PlayTools}@\texttt{7aa1977}：\texttt{PlayLoader.m}，332--334 行 & dylib constructor 自动调用 \texttt{PlayCover.launch}。 \\
\externalmark & 同一 revision：\texttt{PlayCover.swift}，13--23 行 & 运行时初始化顺序及异步启动 MaaTools。 \\
\externalmark & 同一 revision：\texttt{MaaTools.swift}，12--216 行 & 版本、握手、请求帧、listener 和 \texttt{SCRN}。 \\
\externalmark & 同一 revision：\texttt{MaaTools.swift}，284--328 行 & RGB888 转换及 \texttt{BGR\textbackslash{}1} 响应。 \\
\externalmark & 同一 revision：\texttt{AKPlugin.swift}，124--168 行 & ScreenCaptureKit 与旧版 CGWindow 图像采集路径。 \\
\externalmark & 同一 revision：\texttt{PlaySettings.swift}，13--23、76--78 行 & 外部进程读取宿主 plist 及 MaaTools 配置。 \\
\midrule
\multicolumn{3}{p{0.94\textwidth}}{固定 revision 参考：\url{https://github.com/hguandl/PlayTools/tree/7aa19776dbb72b0ae7f0989568f69c7c9c2d6543}。外部源码在当前 PlayCover checkout 中没有以 Carthage/Build 形式存在。} \\
\end{longtable}

\section{验证状态与后续测试矩阵}

本报告已完成宿主侧和固定 revision 外部源码的静态核对。由于当前工作树没有构建好的 PlayTools framework，以下运行时项目尚未宣称通过：

\begin{center}
\begin{tabularx}{\textwidth}{>{\raggedright\arraybackslash}p{0.31\textwidth} X}
\toprule
项目 & 需要的验证 \\
\midrule
协议握手与端序 & 构建目标应用，连接 MaaTools，验证 \texttt{MAA\textbackslash{}0/OKAY}、2 字节大端序和各响应字段。\\
像素格式 & 分别请求 \texttt{SCRN} 与 \texttt{BGR\textbackslash{}1}，检查尺寸、长度、通道顺序和裁剪区域。\\
权限行为 & 在允许和拒绝 Screen Recording 的 TCC 状态下比较日志、空数据和用户反馈。\\
网络边界 & 用 socket 工具确认 listener 的本地/非本地绑定范围，测试端口冲突和重复连接。\\
生命周期 & 测试窗口迟迟不出现、窗口关闭、目标进程退出和 listener 取消。\\
控制命令 & 在受控环境中验证 \texttt{TOUCH} 和 \texttt{TERM} 的授权边界，避免与截图权限混用。\\
\bottomrule
\end{tabularx}
\end{center}

\end{document}
